\scanpage{206}
Considérons le cas où $t'\in S'^+$ est un élément régulier de $S'$; alors $\lambda$ tel que $u(t)=\lambda t'$ est uniquement déterminé par $u$; donc si on se borne aux triples $(S,S^+_{p.d.},t)$ tels que $t$ soit régulier, et aux hom $u:(S,S^+\deleted{,t})\longrightarrow(S',S'^+\deleted{,t'})$ tels que $u(t)\in S.t'$, alors $\mathcal A_*(S,S^+,t)$ devient un foncteur sur cette catégorie, à valeurs ds la cat. des \deleted{\unsure{des}} anneaux gradués \inserted{(\unsure{alternés})} \deleted{\unsure{et}} différentiels filtrés.

La condition envisagée sur $(S,S^+_{pd},t)$ est vérifiée en particulier si on fait $S=k\{T\}$, $S^+=(T)_{pd}$, $t=T$, lorsque $k$ est sans torsion.

Revenons donc à
\[
\mathcal A_*(k)\overset{\mathrm{dfn}}=\mathcal A_*(\underset{k\{T\}}{S_k},S_k^+,T),
\]
et à la détermination des endom de $\mathcal A_*(k)$, comme \deleted{\unsure{anneaux}} \inserted{$k$-algèbre} différentielle $\frac12$ simpliciale (où on ne considère \deleted{\unsure{pas}} la structure de $S$-algèbre, ni la graduation ni les p.d., ni sa filtration a priori-). \deleted{\unsure{Bien à ce propos où}}
Soit $G_k$ le \deleted{groupe} \inserted{monoïde} des couples $(u,\lambda)$, où $u$ est un endom. de $S$ comme $k$-algèbre à p. div., et $\lambda\in S^{\deleted{\unsure{*|+}}}$ tel que $u(T)=\lambda\{T\}T$. Comme $u\deleted{\unsure{(T)}}$ est connu par la connaissance de $u(T)$, $G_k$ est donc le \deleted{groupe} \inserted{monoïde} \unsure{d'ens. sous-jacent} $S$, avec la multiplication
\[
(\lambda'*\lambda)\{T\}\deleted{\unsure{*}}=\lambda'\{T\}\lambda\{T\lambda'\{T\}\}
\]
\deleted{\unsure{de $S^\times$}} --- on désigne par $(u_\lambda)_*$ l'endomorphisme de $\mathcal A_*(k)$ défini par $\lambda\in S$ (et $u_\lambda u_{\lambda'}=u_{\lambda\lambda'}$).
Il est clair que $\lambda\longmapsto u_{\lambda*}$ est un hom de monoïdes \uline{injectif} (en fait $\lambda\longmapsto(u_\lambda)_1$ l'est déjà, car $(u_\lambda)_*(X_0)=\lambda X_0$ et $\mathcal A_1(k)\simeq k\{T\}\{X_0\}$.
Ceci dit,

Si $H_k$ est le monoïde des $\unsure{S_k|G_k}$-endomorphismes de la $k$-algèbre différentielle $\frac12$ simpliciale $\mathcal A_*(k)$, \deleted{et} \deleted{\unsure{est que}} $\Omega_k$ le monoïde de \uline{tous} les endomorphismes de la $k$-alg. diff. $\frac12$ simpliciale, on a
\margin{205}
